% Part: lambda-calculus
% Chapter: syntax
% Section: free-variables

\documentclass[../../../include/open-logic-section]{subfiles}

\begin{document}

\olfileid{lam}{syn}{fv}
\olsection{મુક્ત ચલો}

લૅમ્બડા કલન વિધેયો વિશે છે, અને લૅમ્બડા અમૂર્તીકરણથી વિધેયો
મળે છે. સાહજિક રીતે, $\lambd[x][M]$ એવું વિધેય છે જેમાં વિધેયનો
આર્ગ્યુમેન્ટ~$x$ને સોંપીએ ત્યારે મૂલ્યો~$M$ નક્કી કરે છે. પરંતુ
$M$માં~$x$ની દરેક ઘટના અહીં સંબંધિત નથી: જો~$M$માં બીજું
અમૂર્તીકરણ~$\lambd[x][N]$ હોય તો~$N$માંની~$x$ની ઘટનાઓ
$\lambd[x][N]$ માટે સંબંધિત છે, $\lambd[x][M]$ માટે નહીં. તેથી
$\lambd[x][M]$ની અંદરનું લૅમ્બડા અમૂર્તીકરણ~$\lambd[x]$, $M$માંની
$x$ની જે ઘટનાઓ બીજા કોઈ લૅમ્બડા અમૂર્તીકરણથી અગાઉથી બદ્ધ નથી
તેને \emph{બાંધે} છે---એ~$M$માંની~$x$ની \emph{મુક્ત} ઘટનાઓ છે.

\begin{defn}[વ્યાપ]
પદ~$N$ની અંદર~$\lambd[x][M]$ની ઘટના હોય તો તેમાંની~$M$ની
તેને અનુરૂપ ઘટનાને~$\lambd[x]$નો \emph{વ્યાપ} કહેવાય છે.
\sourcecorrection{OLLAM-011}{સ્થિર અંગ્રેજી મૂળમાં વ્યાપ તરીકે
પદ~$N$ની અનુરૂપ ઘટના લખાઈ છે; વ્યાપ તો અંદરના પદ~$M$ની
અનુરૂપ ઘટના છે. અહીં~$M$ મૂક્યું છે.}
\end{defn}

\begin{defn}[મુક્ત અને બદ્ધ ઘટના]
  પદ~$M$માં ચલ~$x$ની કોઈ ઘટના $\lambd[x]$ના વ્યાપમાં ન હોય તો
  તે \emph{મુક્ત} છે, અને હોય તો \emph{બદ્ધ} છે.
  $\lambd[x][M]$માં ચલ~$x$ની કોઈ ઘટના આરંભના~$\lambd[x]$થી
  ત્યારે અને માત્ર ત્યારે જ બદ્ધ થાય છે જ્યારે~$M$માંની~$x$ની
  તે ઘટના મુક્ત હોય.
\end{defn}

\begin{ex}
  $\lambd[x][x y]$માં $x$ અને~$y$ બંને~$\lambd[x]$ના વ્યાપમાં
  છે; તેથી~$x$ એ~$\lambd[x]$થી બદ્ધ છે. $y$ કોઈ
  $\lambd[y]$ના વ્યાપમાં નથી, તેથી તે મુક્ત છે.
  $\lambd[x][x x]$માં~$x$ની બંને ઘટનાઓ~$\lambd[x]$થી બદ્ધ છે,
  કારણ કે~$xx$માં બંને મુક્ત છે.
  $((\lambd[x][xx])x)$માં~$x$ની છેલ્લી ઘટના મુક્ત છે, કારણ કે
  તે કોઈ~$\lambd[x]$ના વ્યાપમાં નથી.
  $\lambd[x][(\lambd[x][x])x]$માં પહેલા~$\lambd[x]$નો વ્યાપ
  $(\lambd[x][x])x$ છે અને બીજા~$\lambd[x]$નો વ્યાપ
  $x$ની છેલ્લેથી બીજી ઘટના છે. $(\lambd[x][x])x$માં
  $x$ની છેલ્લી ઘટના મુક્ત અને છેલ્લેથી બીજી બદ્ધ છે.
  આમ, $\lambd[x][(\lambd[x][x])x]$માં~$x$ની છેલ્લેથી બીજી
  ઘટના બીજા~$\lambd[x]$થી અને છેલ્લી ઘટના પહેલા
  $\lambd[x]$થી બદ્ધ છે.
\end{ex}

પદ~$P$માં ચલોની બધી ઘટનાઓ તપાસીને તેના મુક્ત ચલોનો ગણ
મેળવી શકાય છે. આ ગણ~$\FV{P}$થી દર્શાવીએ છીએ; તેની સાહજિક
વ્યાખ્યા આ છે:

\begin{defn}[પદના મુક્ત ચલો] \ollabel{def:fv}
  પદના \emph{મુક્ત ચલો}નો ગણ આગમનથી આ રીતે વ્યાખ્યાયિત થાય છે:
  \begin{enumerate}
    \item $\FV{x} = \{x\}$ \ollabel{def:fv1}
    \item $\FV{\lambd[x][N]} = \FV{N} \setminus \{x\}$    \ollabel{def:fv2}
    \item $\FV{PQ} = \FV{P} \cup \FV{Q}$ \ollabel{def:fv3}
  \end{enumerate}
\end{defn}

\begin{prob}
  \begin{enumerate}
  \item આ પદમાં~$\lambd[g]$ અને બે~$\lambd[x]$ના વ્યાપ શોધો:
    $\lambd[g][(\lambd[x][g (x x)]) \lambd[x][g (x x)]]$.
  \item $\lambd[g][(\lambd[x][g (x x)]) \lambd[x][g (x x)]]$માં
    ચલોની બધી ઘટનાઓ બદ્ધ છે? તે અનુક્રમે કયા અમૂર્તીકરણોથી
    બદ્ધ થાય છે?
    \item $\FV{\lambd[x][(\lambd[y][(\lambd[z][x y]) z]) y]}$
    શોધો.
  \end{enumerate}
\end{prob}

\begin{explain}
મુક્ત ચલ બહારના જગત (અર્થાત્~\emph{પરિવેશ}) તરફના નિર્દેશ જેવો
છે. તેથી મુક્ત ચલો ધરાવતું પદ આંશિક રીતે નિર્ધારિત પદ તરીકે
જોઈ શકાય: તેનું વર્તન આપણે પરિવેશ કેવી રીતે ગોઠવીએ તેના પર
આધારિત છે. દાખલા તરીકે, પદ~$\lambd[x][f x]$ આર્ગ્યુમેન્ટ~$x$
સ્વીકારે છે અને તેના પર~$f$ લાગુ કરીને મળતું મૂલ્ય આપે છે.
અહીં ચલ~$f$ મુક્ત છે. પદનું મૂલ્ય તે જે પરિવેશમાં છે તેના પર,
ખાસ કરીને તે પરિવેશમાં~$f$ના મૂલ્ય પર, આધારિત છે.

આ પદ પર અમૂર્તીકરણ લાગુ કરીએ તો
$\lambd[f][\lambd[x][f
    x]]$ મળે છે. હવે આ પદ પરિવેશના ચલ~$f$ પર આધારિત નથી,
કારણ કે તે બે આર્ગ્યુમેન્ટ સ્વીકારીને પહેલાને બીજા પર લાગુ
કરવાનું પરિણામ આપતું વિધેય દર્શાવે છે. પરિવેશમાં~$f$ બદલીશું
તો પણ આ પદના વર્તન પર અસર નહીં થાય: પદ માત્ર આર્ગ્યુમેન્ટ
તરીકે આપેલું મૂલ્ય વાપરશે, પરિવેશમાંના~$f$નું મૂલ્ય નહીં.
\end{explain}

\begin{defn}[બંધ પદ, સંયોજક]
  મુક્ત ચલો વિનાના પદને \emph{બંધ પદ} અથવા
  \emph{સંયોજક} કહે છે.
\end{defn}

\begin{lem}\ollabel{lem:fv}
  \begin{enumerate}
    \item \ollabel{lem:fv-abs} $y \neq x$ હોય તો
      $y \in \FV{\lambd[x][N]}$ ત્યારે અને માત્ર ત્યારે જ
      જ્યારે~$y \in \FV{N}$.
    \item \ollabel{lem:fv-app} $y \in \FV{PQ}$ ત્યારે અને માત્ર
      ત્યારે જ જ્યારે~$y \in \FV{P}$ અથવા~$y
      \in \FV{Q}$.
    \end{enumerate}
\end{lem}

\begin{proof}
  અભ્યાસપ્રશ્ન.
\end{proof}

\begin{prob}
  \olref[lam][syn][fv]{lem:fv} સાબિત કરો.
\end{prob}

\end{document}
